% Данный файл распространяется под лицензией CC BY 4.0.
% Текст лицензии: https://creativecommons.org/licenses/by/4.0/
% Компиляция: XeLaTeX

\documentclass[a4paper]{article}

\usepackage[a4paper, top=10mm, bottom=12mm, left=12mm, right=12mm]{geometry}
\usepackage{polyglossia}
\setdefaultlanguage[babelshorthands=true]{russian}
\setotherlanguage{english}
\usepackage{fontspec}
\setmainfont{FreeSerif}
\setmonofont[Scale=0.82]{DejaVu Sans Mono}
\usepackage[tiny, compact]{titlesec}
\usepackage{titling}
\setlength{\droptitle}{-1.2cm}
\pretitle{\begin{center}\begin{bfseries}\Large}
\posttitle{\par\end{bfseries}\end{center}}
\preauthor{\begin{center}\normalsize}
\postauthor{\par\end{center}\vspace{-1.6cm}}
\usepackage{hyperref}
\usepackage{bookmark}
\usepackage{csquotes}
\hypersetup{colorlinks=true, linkcolor=black, citecolor=black, urlcolor=blue}
\urlstyle{same}

\title{Бенчмаркинг обработки скетчей в SBC MarLine}
\author{Шахов Андрей}
\date{}

\begin{document}
\maketitle

\begin{flushright}
    Группа: \emph{25.Б72-мм}\\
    Кафедра: \emph{системного программирования}\\
    Научный руководитель: \emph{Гориховский Вячеслав Игоревич}\\
    Консультант: \emph{Дмитриевцев Алексей Сергеевич, младший инженер по разработке ПО, ООО "КНС ГРУПП"}\\
    Номер семестра практики: \emph{3}
\end{flushright}

\section{Введение}

\section{Постановка задачи}

\textbf{Цель работы} --- разработать воспроизводимую методику бенчмаркинга обработки скетчей в SBC MarLine для оценки стоимости вставки и поиска в различных реализациях индекса при разных размерах базы, а также спроектировать общий интерфейс измерения задержек, применимый ко всем этапам конвейера.

\textbf{Для достижения цели требуется решить следующие задачи:}
\begin{enumerate}
    \item Определить стадии конвейера, границы отдельных измерений, измеряемые величины и правила агрегации результатов.
    \item Подготовить воспроизводимую тестовую нагрузку на архивах версий ядра Linux и зафиксировать параметры построения скетчей.
    \item Разработать и проверить сценарии измерения вставки, поиска top-$k$, полного прохода по запросам и построения индекса.
    \item Сравнить линейную, обратную и эвристическую реализации индекса при разных размерах базы; для эвристического поиска оценить качество результатов относительно полного поиска.
    \item Спроектировать потокобезопасный интерфейс сбора задержек, пригодный для последовательной и параллельной обработки.
    \item Предусмотреть агрегацию статистик и экспорт результатов в машиночитаемом формате.
    \item Подготовить тесты и провести экспериментальную оценку точности и накладных расходов измерительного инструмента.
\end{enumerate}

\section{Обзор}

\section{Описание решения}

\section{Апробация и эксперименты}

\section{Заключение}

Материалы практики и исходный код: \url{https://github.com/kevetka/marline_benchmarks}.

\begin{thebibliography}{9}
    \bibitem{aronovich2009}
    Aronovich L., Asher R., Bachmat E., Bitner H., Hirsch M., Klein S.~T.
    \newblock The Design of a Similarity Based Deduplication System.
    \newblock In: Proceedings of SYSTOR 2009. 2009. P.~1--14.
    \newblock DOI: \href{https://doi.org/10.1145/1534530.1534539}{10.1145/1534530.1534539}.

    \bibitem{finesse}
    Zhang Y., Xia W., Feng D., Jiang H., Hua Y., Wang Q.
    \newblock Finesse: Fine-Grained Feature Locality Based Fast Resemblance Detection for Post-Deduplication Delta Compression.
    \newblock In: Proceedings of FAST 2019. 2019. P.~121--128.
    \newblock URL: \url{https://www.usenix.org/conference/fast19/presentation/zhang}.

    \bibitem{odess}
    Zou X., Deng C., Xia W., Shilane P., Tan H., Zhang H., Wang X.
    \newblock ODESS: Speeding up Resemblance Detection for Redundancy Elimination by Fast Content-Defined Sampling.
    \newblock In: Proceedings of ICDE 2021. 2021. P.~480--491.
    \newblock DOI: \href{https://doi.org/10.1109/ICDE51399.2021.00048}{10.1109/ICDE51399.2021.00048}.

    \bibitem{fastcdc}
    Xia W., Zou X., Jiang H., Zhou Y., Liu C., Feng D., Hua Y., Hu Y., Zhang Y.
    \newblock The Design of Fast Content-Defined Chunking for Data Deduplication-Based Storage Systems.
    \newblock IEEE Transactions on Parallel and Distributed Systems. 2020. Vol.~31, no.~9. P.~2017--2031.
    \newblock DOI: \href{https://doi.org/10.1109/TPDS.2020.2984632}{10.1109/TPDS.2020.2984632}.

    \bibitem{inverted}
    Zobel J., Moffat A.
    \newblock Inverted Files for Text Search Engines.
    \newblock ACM Computing Surveys. 2006. Vol.~38, no.~2. Article~6.
    \newblock DOI: \href{https://doi.org/10.1145/1132956.1132959}{10.1145/1132956.1132959}.
\end{thebibliography}

\end{document}
